\chapter{Production Backreaction and the Frozen Benchmark}
\label{ch:backreaction}

\section{Control conditions for the rapid-crossing spectrum}

The Gaussian occupation \eqref{eq:univ-distribution} follows from a prescribed approximately linear crossing.  Three independent small parameters are required near the populated band:
\begin{equation}
\epsilon_H
:=\frac{H_\star^2}{A_\star m_{\mathrm{peak}}}\ll1,
\label{eq:back-hubble-control}
\end{equation}
\begin{equation}
\epsilon_{\mathrm{lin}}
:=\frac{|\ddot\phi_\star|}
{|\dot\phi_\star|\sqrt{A_\star m_{\mathrm{peak}}}}\ll1,
\label{eq:back-linear-control}
\end{equation}
\begin{equation}
\epsilon_{\mathrm{br}}
:=\frac{\widetilde\rho_\star}{K_\star}\ll1.
\label{eq:back-br-control}
\end{equation}
The first two control expansion and curvature across the event.  The third controls depletion by the produced modes.  A posteriori normalization cannot replace \eqref{eq:back-br-control} because depletion changes the background that generated the Bogoliubov coefficients.

\section{Benchmark energy-overdraw theorem}

The benchmark with $\lambda=1$ defines a shape-preserving projection
\begin{equation}
f_h(k)\longmapsto\eta f_h(k),
\qquad
\eta=\min\left(1,\frac{K_\star}{\widetilde\rho_\star}\right).
\label{eq:back-projection}
\end{equation}

\begin{theorem}[Failure of negligible backreaction]
For the specified benchmark record,
\begin{equation}
\eta=\BridgeEta
\end{equation}
and therefore
\begin{equation}
\boxed{
\epsilon_{\mathrm{br}}
=\frac1\eta
=\BridgeOverdraw.}
\label{eq:back-overdraw}
\end{equation}
The benchmark is outside the negligible-backreaction domain by more than twenty orders of magnitude.
\end{theorem}

\begin{proof}
The value of $\eta$ is read from the frozen discrete-tower calculation.  On the active branch of \eqref{eq:back-projection}, $\eta=K_\star/\widetilde\rho_\star$.  Inversion gives \eqref{eq:back-overdraw}.
\end{proof}

\begin{negativebox}[Physical consequence]
The exact ratios $w_\star=1/9$ and $f_R=525\pi/2048$ are properties of the projected continuous ansatz at $\sigma=1$.  They are not predictions of the coupled crossing until a backreacting solution reproduces the same spectral shape.  With $\epsilon_{\mathrm{br}}\gg1$, strong depletion is mandatory.
\end{negativebox}

\section{Hubble-loss correction before the crossing}

If the average field begins at rest at $\phi_i$ and reaches $\phi=0$, the mechanical crossing energy in FLRW obeys
\begin{equation}
\boxed{
K_{\star^-}
=U(\phi_i)-U(0)
-3\int_{t_i}^{t_\star}H\dot\phi^2\,\dd t
-E_{\mathrm{other}},}
\label{eq:back-hubble-loss}
\end{equation}
where $E_{\mathrm{other}}$ denotes the integrated work removed from the average field before the crossing.  It need not equal the final redshifted energy density of the receiving sector.  The formula without the integral is exact only in a negligible-expansion and negligible-transfer limit.  The benchmark does not fix the absolute ratio $\mu/H_\star$, so that approximation cannot be checked from the dimensionless record alone.

\section{Continuous order-one backreaction estimate}

In the arbitrary-order continuous model,
\begin{equation}
\widetilde\rho_\star
=C_E(\sigma)\frac{A^{\alpha+4}}{\Lambda^\alpha},
\qquad
K_\star=\frac{v^2A^2}{2\gamma},
\qquad \gamma>0.
\end{equation}
Introduce
\begin{equation}
a=\frac A\mu,
\qquad
\ell=\frac\Lambda\mu,
\qquad
\frac{v^2}{\mu^2}=\frac1\lambda.
\end{equation}
The raw continuous overdraw is
\begin{equation}
R_{\mathrm c}(a)
=2\gamma\lambda C_E(\sigma)
\frac{a^{\alpha+2}}{\ell^\alpha}.
\end{equation}
For the frozen dimensionless input, $R_{\mathrm c,in}=\BridgeContinuumOverdraw$.  If changing $A$ were the only response, the order-one boundary $R_{\mathrm c}=1$ would satisfy
\begin{equation}
\boxed{
\frac{A_{\mathrm{br,c}}}{A_{\mathrm{in}}}
=R_{\mathrm c,in}^{-1/(\alpha+2)}.}
\label{eq:back-cont-slow}
\end{equation}
For the benchmark, $\alpha=1$ and
\begin{align}
\frac{A_{\mathrm{br,c}}}{A_{\mathrm{in}}}
&=\BridgeContinuumSlow,\\
\frac{A_{\mathrm{br,c}}}{\mu}
&=\BridgeContinuumA,\\
\frac{\Lambda}{A_{\mathrm{br,c}}}
&=\BridgeContinuumX.
\end{align}
The lightest level would then have suppression exponent
\begin{equation}
\pi\frac{m_1}{A_{\mathrm{br,c}}}
\simeq\BridgeContinuumExponent.
\label{eq:back-cont-exponent}
\end{equation}
The continuous populated-band approximation is self-inconsistent at that point.  This is not a paradox.  It says that slowing the same continuous spectrum cannot reach its order-one boundary while holding the microscopic tower fixed.  The raw continuous coefficient is used here rather than the fixed-prime projection factor $\eta$.

\section{Direct discrete order-one threshold}

The next calculation keeps the benchmark's fixed $\Lambda$, $p$, $\sigma$, $\gamma$, $v$, and Gaussian single-crossing formula.  It changes only $A$ and evaluates the cutoff-controlled discrete momentum moments.  Define
\begin{equation}
R_{\mathrm{disc}}(A)
:=\frac{\widetilde\rho_{\mathrm{disc}}(A)}{v^2A^2/(2\gamma)}.
\label{eq:back-Rdisc}
\end{equation}

\begin{diagnostic}[Discrete backreaction threshold]
The equation $R_{\mathrm{disc}}(A)=1$ has the numerical solution
\begin{equation}
\boxed{
\frac{A_{\mathrm{br,d}}}{\mu}=\BridgeDiscreteA,
\qquad
\frac{\Lambda}{A_{\mathrm{br,d}}}=\BridgeDiscreteX.}
\label{eq:back-disc-sat}
\end{equation}
At this point,
\begin{equation}
\boxed{
w_{\mathrm{br,d}}=\BridgeDiscreteW,
\qquad
f_{R,\mathrm{br,d}}=\BridgeDiscreteRest,
\qquad
\frac{\langle m\rangle_R}{A}=\BridgeDiscreteMeanMass.}
\label{eq:back-disc-ratios}
\end{equation}
\end{diagnostic}

The calculation uses 128-node generalized Gauss--Laguerre momentum quadrature and the exact discrete degeneracies.  Three levels fall below the declared suppression cutoff.  The first carries essentially all of the production weight.

\begin{figure}[htbp]
\centering
\includegraphics[width=0.98\linewidth]{backreaction_obstruction.pdf}
\caption{Left: the unprojected energy relative to the crossing reservoir as $\Lambda/A$ is increased at fixed microscopic $\Lambda$.  The order-one backreaction threshold occurs after the continuous populated-band regime has been left.  Right: the production state becomes more rest dominated and its pressure falls as the tower sparsifies.}
\label{fig:backreaction-obstruction}
\end{figure}

\begin{statusbox}[Diagnostic boundary]
Equation \eqref{eq:back-disc-sat} is not a time-dependent backreaction solution and is not an energy-partition equation for the original crossing.  It marks where the prescribed spectrum and kinetic reservoir become comparable inside a one-parameter Gaussian family.  Its value is useful because the family has already become sparse and cold at that boundary.  A dynamical solution can depart further.
\end{statusbox}

\section{Interaction-strength obstruction}

For a detail eigenvalue $\lambda_h=\Lambda^2p^{2\sigma h}$, the canonical four-dimensional action contains the quartic interaction
\begin{equation}
-\frac12g_h\phi^2\chi_h^2,
\qquad
g_h=\gamma\frac{\lambda_h}{v^2}.
\label{eq:back-gh}
\end{equation}
Using $v^2=\mu^2/\lambda$ and $m_h^2\simeq\lambda_h$ gives
\begin{equation}
g_h\simeq\gamma\lambda\left(\frac{m_h}{\mu}\right)^2.
\label{eq:back-gh-scaled}
\end{equation}
At the frozen discrete peak with $\gamma=16$ and $\lambda=1$,
\begin{equation}
\boxed{g_{\mathrm{peak}}\simeq\BridgePeakCoupling.}
\label{eq:back-peak-coupling}
\end{equation}
A nominal loop-control condition $g_{\mathrm{peak}}<1$ would instead require
\begin{equation}
\lambda<\BridgeLambdaBound
\label{eq:back-lambda-bound}
\end{equation}
with the other benchmark ratios held fixed.  At $\lambda=1$, even the lightest level at the discrete threshold has
\begin{equation}
g_1\simeq\BridgeLowestCoupling.
\end{equation}

At fixed calibrated ratios $A/\mu$ and $\Lambda/\mu$, the two diagnostics are parametrically linked:
\begin{equation}
R_{\mathrm{br}}(\lambda)
=\frac{\lambda}{\lambda_{\mathrm{frozen}}}
R_{\mathrm{br}}(\lambda_{\mathrm{frozen}}).
\end{equation}
For $\lambda_{\mathrm{frozen}}=1$, avoiding energy overdraw requires
\begin{equation}
\boxed{\lambda\le\BridgeEnergyLambdaBound.}
\end{equation}
At that boundary, the peak and lightest-level couplings would be
\begin{equation}
g_{\mathrm{peak}}=\BridgePeakAtEnergyBound,
\qquad
g_1=\BridgeLowestAtEnergyBound.
\end{equation}

\begin{negativebox}[Perturbative boundary]
The printed $\lambda=1$ benchmark is neither weakly backreacting nor under loop control.  Large $g_h$ does not by itself invalidate the exact solution of a quadratic mode equation on a prescribed background.  The fixed-background approximation fails independently because of the energy overdraw.  Reducing the free quartic can cure both tree-level diagnostics along this parameter slice, but it changes the absolute scale $\mu\propto\lambda^{1/4}$ at fixed vacuum density and requires new Hubble, production, and loop checks.
\end{negativebox}

\section{Self-consistent semiclassical system}

A minimal next calculation must solve the background, modes, and geometry together.  At finite depth, one suitable Hartree system is
\begin{align}
\ddot\phi+3H\dot\phi+U'(\phi)
+\frac12F'(\phi)Y_H^{\mathrm{ren}}&=0,
\label{eq:back-hartree-phi}\\
\ddot u_{h,k}+3H\dot u_{h,k}
+\left(\frac{k^2}{a^2}+m_b^2+F(\phi)\lambda_h\right)u_{h,k}&=0,
\label{eq:back-hartree-mode}\\
3M_{\mathrm{Pl}}^2H^2
&=\frac12\dot\phi^2+U(\phi)+\rho_H^{\mathrm{ren}}+\rho_{\mathrm{vis}},
\label{eq:back-hartree-friedmann}
\end{align}
where
\begin{equation}
Y_H^{\mathrm{ren}}
=\sum_{h=1}^{H}d_h\lambda_h
\int\frac{\dd^3k}{(2\pi)^3}
\left(|u_{h,k}|^2-\text{adiabatic subtraction}\right).
\label{eq:back-Yren}
\end{equation}
The same subtraction prescription must be used in $\rho_H^{\mathrm{ren}}$.  Equations \eqref{eq:cov-cont-chi}--\eqref{eq:cov-cont-phi} then provide an exact numerical conservation test.

The calculation is not complete until it demonstrates all of the following.
\begin{enumerate}[leftmargin=2.2em]
\item cutoff stability over a declared range of $H$ and momentum resolution
\item conservation after renormalization
\item independence from an arbitrary scale-factor normalization
\item a spectral output rather than a shape-preserving projection
\item a regime in which quantum and higher-order corrections are controlled
\end{enumerate}

No result in this chapter claims that such a solution exists for the frozen parameters.  The energy and interaction-strength diagnostics make survival of that exact benchmark unlikely without a substantial microscopic revision.
